
\subsection{Model Presentation}
\label{sec:results_model_presentation}

The evaluated models operate on the event representation defined in
Section~\ref{sec:method_features} and use the canonical-\(y\) target
convention of Section~\ref{sec:method_targets}. Their common primary
prediction is the dominant-electron assignment of every retained
\gls{ecal} hit. This establishes a direct measure of shower separation before
considering the additional event-level outputs of the extended model.

\subsubsection{Controlled Hit-Classification Comparisons}

The baseline comparison consists of four models with one shared training
target. The \textit{ECalGravNet} and \textit{ECalTransformer} models receive
only reconstructed \gls{ecal} hit position and energy. They test how far
overlapping showers can be separated using the calorimeter pattern itself.
The \textit{ECalTpadGravNet} and \textit{ECalTpadTransformer} models receive
the corresponding \gls{ecal} hits together with trigger-pad-track centroid
and photoelectron tokens. Comparing ECal-only and context-aware variants
therefore isolates the value of the additional detector information without
changing the hit-level prediction task.

The GravNet models build latent neighbourhoods of detector objects and use
local feature propagation before predicting an origin class for each ECal
hit. The transformer models instead apply full self-attention over all input
objects in an event. Thus, the former expresses a learned local graph
inductive bias, while the latter allows direct global conditioning across
potentially overlapping showers. Both are appropriate for an unordered,
variable-size detector response and can incorporate trigger-pad-track
context without imposing an image representation.

\subsubsection{Multi-Task Context Model}

The MLPF-lite transformer uses the combined ECal and trigger-pad-track input
and supplements dominant-origin classification with a per-hit prediction of
the energy fraction contributed by each electron. It is included as an
intermediate multi-task reconstruction model: the dominant class indicates
the assigned shower, whereas the fraction output describes ambiguity in
overlap regions. Although inspired by \gls{mlpf}, this output is deliberately
limited to the LDMX ECal shower-association task rather than a complete
particle-flow reconstruction.

\subsubsection{Slot and Electron-Count Model}

For mixed two-electron and three-electron samples, the principal extended
model is the ECal-plus-trigger-pad slot transformer. Its per-hit classes
represent background/noise and up to three spatially ordered electron slots.
In addition to hit-origin and energy-fraction heads, it uses a pooled event
representation to predict which electron slots are present and the total
electron multiplicity. This model is the natural endpoint of the study:
shower separation is learned at hit level, while the event-level head tests
whether the separated activity supports identification of an additional
overlapping electron. The following sections report training behaviour and
the corresponding hit-level and event-level diagnostic outputs without
altering this model definition.

\subsection{Training Performance}
Present loss curves.
Discuss convergence and stability.

*Figure with loss curve and more*

*Table with results in numbers*

\subsection{Shower Separation Performance}
Compare GNN with baseline methods.
Evaluate electron counting efficiency.
Analyze difficult pile-up scenarios.

*Figure with histograms containing results of GNN vs baseline*

*Table with results in numbers*

\subsection{Energy Reconstruction Quality}
Discuss energy resolution and purity.
Present quantitative comparisons.

*Figure with histograms containing results of GNN vs baseline*

*Table with results in numbers*

\subsection{Robustness Studies}
Evaluate performance under varying pile-up conditions.
Study sensitivity to noise or detector effects.

*Figure with histograms containing results of GNN vs baseline*

*Table with results in numbers*

\subsection{More...?}
